% ...existing code do resumo...
Este trabalho apresenta o desenvolvimento e análise de dois sistemas fuzzy: um sistema Mamdani para controle de ventilação em ambientes fechados e um sistema Takagi-Sugeno para aproximação de funções não lineares. O sistema Mamdani foi projetado para sugerir a intensidade de ventilação a partir de variáveis ambientais e ocupacionais, utilizando funções de pertinência triangulares e diferentes técnicas de defuzzificação. O sistema Takagi-Sugeno foi implementado para aproximar a função $f(x) = e^{-x/5} \cdot \sin(3x) + 0.5 \cdot \sin(x)$, avaliando diferentes funções de pertinência, operadores e métodos de otimização dos consequentes. Os resultados demonstram que ambos os sistemas são eficazes em suas propostas, com o Mamdani fornecendo decisões qualitativas interpretáveis e o Takagi-Sugeno apresentando boa precisão na aproximação da função alvo, especialmente após otimização dos parâmetros. O estudo evidencia a importância da escolha adequada das funções de pertinência, operadores e técnicas de otimização para o desempenho dos sistemas fuzzy.

\bigskip

\textbf{Abstract} \\
This work presents the development and analysis of two fuzzy systems: a Mamdani system for ventilation control in indoor environments and a Takagi-Sugeno system for nonlinear function approximation. The Mamdani system was designed to suggest ventilation intensity based on environmental and occupancy variables, using triangular membership functions and different defuzzification techniques. The Takagi-Sugeno system was implemented to approximate the function $f(x) = e^{-x/5} \cdot \sin(3x) + 0.5 \cdot \sin(x)$, evaluating different membership functions, operators, and consequent optimization methods. The results show that both systems are effective in their proposals, with the Mamdani system providing interpretable qualitative decisions and the Takagi-Sugeno system achieving good accuracy in approximating the target function, especially after parameter optimization. The study highlights the importance of proper selection of membership functions, operators, and optimization techniques for the performance of fuzzy systems.
